\section{讨论}

\subsection{限定与推广}

本章所发展的方法，促进从定性因果假设（编码于图中）与非实验观测的结合中得出定量的因果推断。这些因果假设本身通常无法在非实验研究中被检验，除非它们对被观测分布施加约束。最常见的一类约束以条件独立的形式出现，正如各图中的 d-分离条件所传达的那样。另一类约束取数值不等式的形式。例如在第八章我们将表明，与工具变量相联系的假设（图 3.7(b)）须接受条件概率不等式形式的否证检验（Pearl 1995b）。尽管如此，这类约束只能检验图中体现的因果假设的一小部分；这些假设的大部分必须由领域知识来证实，而领域知识或来自理论考量（例如“气压计下降不会导致下雨”），或来自相关的实验研究。例如，Moertel 等 (1985) 否证“维生素 C 对癌症有效”这一假说的实验研究，可以作为涉及维生素 C 与癌症病人的观测研究中的实质性假设；在相应的图中，它将被表示为（维生素 C 与癌症之间的）一条缺失连线。概括而言，本章所述方法的主要用途不在于检验因果假设，而在于提供一种有效的语言，使这些假设变得精确而明确。假设由此可以被分离出来，以供审议或实验；一旦得到验证，便可与统计数据结合，产生因果效应的定量估计。

本书只简要讨论（见 \S 8.5）的一个重要问题是抽样变异性。因果效应估计量的数学推导，应被视为朝着“为这些估计量补充置信区间与显著性水平”迈出的第一步，如同受控实验的传统分析那样。不过我们应当指出，获得因果效应的非参数估计量，并不意味着在研究估计阶段应当弃用参数形式。例如，若高斯扰动、零均值、可加交互这些假设被认为是合理的，则 \textup{(3.28)} 所给的估计量可化为乘积 $E(Y\mid\hat{x})=r_{ZX}\,r_{YZ\cdot X}\,x$，其中 $r_{YZ\cdot X}$ 是标准化回归系数（\S 5.3.1）；估计问题于是归结为估计回归系数（例如用最小二乘）。更精细的估计技术可见 Rosenbaum and Rubin (1983)、Robins (1989, sec.~17) 与 Robins 等 (1992, pp.~331--3)。例如，Rosenbaum and Rubin (1983) 的“倾向评分”方法在被调整协变量的维数很高时颇为有用（\S 11.3.5）。Robins (1999) 表明，与其估计 \textup{(3.19)} 调整公式中的各个因子，往往更有利的是使用
\[
P(y\mid\hat{x})=\sum_{z}\frac{P(x,y,z)}{P(x\mid z)},
\]
其中干预前分布保持未分解。于是可以单独估计分母 $P(x\mid z)$，以该估计的倒数对各个样本加权，并把加权后的样本当作从干预后分布 $P(y\mid\hat{x})$ 中随机抽取的一样处理。干预后参数（如 $E(Y\mid\hat{x})$）便可用普通最小二乘估计。这一方法在具有时变协变量的纵向研究中尤为有利，如 \S 3.2.3（见 \textup{(3.18)}）与 \S 4.4.3 所讨论的问题。

本章所提方法有若干值得注意的推广。第一，原子干预的识别分析可以推广到复杂的时间变化政策：其中受控变量集 $X$ 通过函数式或随机式策略，以规定方式响应某组协变量 $Z$，如 \S 3.2.3 与 \S 4.4.3 所示。第四章（\S 4.4.3）将表明，识别此类政策的效应需要在相应图中满足一系列后门条件。

第二项推广涉及在非递归模型、即含指向环或反馈环的因果图中使用干预演算（定理 3.4.1）。以“从模型中抹去方程”定义因果效应的基本做法（定义 3.2.1）仍可沿用于非递归系统（Strotz and Wold 1960; Sobel 1990），但此时须处理两个问题。第一，识别分析必须确保剩余子模型的稳定性（Fisher 1970）。第二，DAG 的 d-分离准则必须推广到循环图。d-分离的有效性已在非递归线性模型（Spirtes 1995）以及涉及离散变量的非线性系统（Pearl and Dechter 1996）中得到确立。然而，在循环非线性系统中，因果效应估计量的计算将更为困难，因为把 $P(y\mid\hat{x})$ 符号化简为无帽表达式可能需要解非线性方程。第七章（\S 7.2.1）演示了非递归线性系统中政策与反事实的求值（另见 Balke and Pearl 1995a）。

第三项推广涉及把干预演算（定理 3.4.1）推广到数据并非在独立同分布（i.i.d.）抽样下获得的情形。例如可以设想，一位医生仅在既往病人存活比例降至某阈值以下时，才给病人开出某种治疗。此类情形要求从非独立样本中估计因果效应 $P(y\mid\hat{x})$。Vladimir Vovk (1996) 给出了在抽样非 i.i.d. 时定理 3.4.1 各规则适用的条件，并进一步把三条推理规则表述为一个逻辑产生式系统。

\subsection{图作为数学语言}

把实质性背景知识纳入概率推断的益处，早在 Thomas Bayes (1763) 与 Pierre Laplace (1814) 时便已被认识；它在复杂统计研究的分析与解释中的关键作用，也得到多数现代统计学家的普遍承认。然而，可用于表达背景知识的数学语言，一直处于相当可怜的欠发达状态。

传统上，统计学家只认可一种把实质知识与统计数据结合的方式：即给分布参数赋予主观先验的贝叶斯方法。为在这一框架内纳入因果信息，诸如“$Y$ 不受 $X$ 影响”这样的朴素因果陈述，必须被转换为能够接受概率值的语句或事件（例如反事实）。比如，为传达“泥不会导致雨”这一无害的假设，我们不得不使用一个相当不自然的表达，说反事实事件“若当时没有泥，则会下雨”的概率，等于“若当时有泥，则会下雨”的概率。确实，Neyman–Rubin 的潜在结局进路正是以这种方式获得统计上的正当性：因果判断被表述为涉及反事实变量之概率函数的约束（见 \S 3.6.3）。

因果图提供了另一种把数据与因果信息结合的语言。这种语言把朴素因果陈述接受为基本原初概念，从而简化了贝叶斯路线。此类陈述仅仅指明两个所关心变量之间是否存在因果联系，它们在日常话语中被普遍使用，为科学家交流经验、组织知识提供了自然的方式。\footnote{值得注意的是，本章的许多读者（包括本书的两位审稿人）把这里所述的方法归入“贝叶斯阵营”，并认为它依赖于一个“好的先验”。这种归类是误导性的。该方法确实依赖于主观假设（如泥不会导致雨），但这类假设是\emph{因果}的，而非统计的，无法被表述为联合分布参数上的先验概率。}因此可以预期，因果图的语言将在需要大量领域知识的问题中找到应用。

这种语言并不新鲜。用图与结构方程模型传达因果信息，在社会科学与计量经济学中已相当流行。然而统计学家通常对这些模型心存疑虑，也许是因为社会科学家与计量经济学家未能为其模型的经验内容提供无歧义的定义——即未能指明：在何种实验条件下（哪怕是假设性的），给定的结构方程将约束其结果。（第五章讨论结构方程在社会科学与经济学中的历史。）结果，甚至连“结构系数”或“缺失连线”这样基本的概念，都成为激烈争议（Freedman 1987; Goldberger 1992）与误读（Whittaker 1990, p.~302; Wermuth 1992; Cox and Wermuth 1993）的对象。

这一争议与误解的历史，在很大程度上源于缺乏一套恰当的数学记号来定义因果建模的基本概念。例如，标准概率记号无法表达结构方程 $y=bx+e_Y$ 中系数 $b$ 的经验内容，即便人们准备假定 $\sigma_{XY}$（一个未观测量）与 $X$ 不相关。分析者把不“直接影响”$Y$ 的变量排除在方程之外，这一做法也无法赋予任何概率含义。\footnote{Goldberger (1992) 与 Wermuth (1992) 之间的激烈争论，正围绕 Wermuth 坚持要给结构方程中的零系数赋予统计解释这一点展开（见 \S 5.4.1）。}

本章所发展的记号赋予这些（因果）概念以清晰的经验解释，因为它使人能够精确指明：在给定实验中，什么是被保持恒定的，什么只是被测量的。（这一区分的必要性已为许多研究者所认识，其中最著名的是 Pratt and Schlaifer 1988 与 Cox 1992。）$b$ 的含义就是 $\partial E(Y\mid\hat{x})/\partial x$，即在 $X$ 被外部控制保持在 $x$ 的实验中，$Y$ 之期望（随 $x$）的变化率。这一解释成立与否，与 $e_Y$ 和 $X$ 是否相关（例如经由另一个方程 $x=ay+e_x$）无关。同样，分析者关于哪些变量应纳入给定方程的决断，可以立足于一个假设性的受控实验：若在其余一切变量（$S\backslash YZ$）保持恒定时 $Z$ 对 $Y$ 没有影响（对每个 $e_Y$ 水平皆然），则把 $Z$ 从 $Y$ 的方程中排除；这意味着 $P(y\mid\hat{z},\hat{s}_{YZ})=P(y\mid\hat{s}_{YZ})$。具体地说，从方程 $y=bx+e_Y$ 中排除的变量，并不是在给定 $X$ 的测量时与 $Y$ 条件独立，而是在给定 $X$ 的设定时与 $Y$ 因果无关。“扰动项”$e_Y$ 的操作含义同样被祛魅：$e_Y$ 被定义为差值 $Y-E(Y\mid\hat{x})$。两个扰动项 $e_X$ 与 $e_Y$ 相关，当且仅当 $P(y\mid\hat{x},\hat{s}_{XY})\ne P(y\mid x,\hat{s}_{XY})$，如此等等（进一步的阐述见第五章 \S 5.4）。

帽记号所提供的区分澄清了结构方程的经验基础，应当能使因果模型更易为经验研究者接受。此外，由于大多数科学知识是围绕“把 $X$ 保持固定”而非“对 $X$ 取条件”这一操作组织起来的，本章所发展的记号与演算，应当能为科学家提供一种有效手段，以交流实质信息并推断其逻辑后果。

\subsection{从图到潜在结局的翻译}

本章使用两种因果信息的表示：图与结构方程，前者是后者的抽象。两种表示在近一个世纪里都备受争议。一方面，经济学家与社会科学研究者接纳了这些建模工具，却持续质疑并争论他们所估计参数中的因果内容（详见 \S 5.1 与 \S 5.4）；结果，结构模型在政策制定语境中的使用常被视为可疑。统计学家大体上拒斥这两种表示，认为它们有问题（Freedman 1987）乃至毫无意义（Wermuth 1992; Holland 1995），有时在被迫交流因果信息时转而求助于 Neyman–Rubin 潜在结局记号（Rubin 1990）。\footnote{计量经济学文献中在“切换回归”名目下发展了一个平行框架（Manski 1995, p.~38），Heckman (1996) 将其归功于 Roy (1951) 与 Quandt (1958)；但由于缺乏 \textup{(3.51)} 那样的形式语义，它并未超出“框架”的阶段。}关于结构进路与潜在结局进路之间关系的详细形式分析，见第七章（\S 7.4.4），那里证明了二者的数学等价性——一个进路中的定理蕴含另一个进路中的定理。本节着重指出关键的方法论差异。

潜在结局框架中分析的原初对象是基于个体的响应变量，记作 $Y(x,u)$ 或 $Y_x(u)$，读作：“在个体 $u$ 中、若 $X$ 曾为 $x$，$Y$ 将会取得的值”。这一反事实表达式在结构方程模型中有形式解释。考虑一个含方程组
\begin{equation}
x_i=f_i(\pa_i,u_i),\qquad i=1,\dots,n, \tag{3.50}
\end{equation}
的结构模型 $M$，如 \textup{(3.4)} 那样。设 $U$ 表示背景变量向量 $(U_1,\dots,U_n)$，$X$ 与 $Y$ 为被观测变量的两个不相交子集，$M_x$ 为按定义 3.2.1 把 $X$ 中变量所对应方程替换为 $X=x$ 所得的子模型。$Y(x,u)$ 的结构解释由
\begin{equation}
Y(x,u)\triangleq Y_{M_x}(u) \tag{3.51}
\end{equation}
给出。也就是说，$Y(x,u)$ 是 $M$ 的子模型 $M_x$ 中、在实现 $U=u$ 之下 $Y$ 的（唯一）解。尽管潜在结局文献中的“个体”一词通常指总体中某个具体个人的身份，个体也可以被设想为刻画该个人的一组属性、所研究的实验条件、一天中的时刻等等——所有这些在结构建模中都表示为向量 $u$ 的分量。事实上，对 $U$ 的唯一要求是：(i) 它表示足够多的背景因素，使内生变量之间的关系成为确定性的；(ii) 数据由从 $P(u)$ 抽取的独立样本构成。实验中某个具体个人的身份往往足以满足这一要求，因为它代表该个体的解剖与遗传构造，而这些通常足以决定该个体对处理或其他所关心项目的反应。

式 \textup{(3.51)} 定义了那个晦涩的英语短语“在个体 $u$ 中、若 $X$ 曾为 $x$ 时 $Y$ 会取得的值”与“把 $X$ 的变化转化为 $Y$ 之变化的物理过程”之间的关键形式联系。子模型 $M_x$ 的构造，精确阐明了假设性短语“若 $X$ 曾为 $x$”可以如何被实现，以及必须让什么过程作出让步才能使 $X=x$ 成为现实。

有了对 $Y(x,u)$ 的这一解释，对比潜在结局框架与结构框架中的因果推断方法论便很有教益。若把 $U$ 视为随机变量，则反事实 $Y(x,u)$ 的值也成为随机变量，记作 $Y(x)$ 或 $Y_x$。潜在结局分析的做法是：把被观测分布 $P(x_1,\dots,x_n)$ 想象为定义在被观测变量与反事实变量二者之上的增广概率函数 $P^*$ 的边缘分布。关于因果效应的查询（在我们的结构分析中写作 $P(y\mid\hat{x})$）被表述为关于所关心反事实变量之边缘分布的查询，写作 $P^*(Y(x)=y)$。这些新的假设性实体 $Y(x)$ 被当作普通随机变量；例如，假定它们服从概率演算的公理、条件化定律与条件独立公理。此外，这些假设性实体被假定通过一致性约束（Robins 1986）与可观测变量相联系，例如\footnote{Gibbard and Harper (1976, p.~156) 把这一约束表述为 $A\Rightarrow[(A\ \square\!\!\to S)\equiv S]$。}
\begin{equation}
X=x\Rightarrow Y(x)=Y, \tag{3.52}
\end{equation}
它说的是：对每个 $u$，若 $X$ 的实际值恰为 $x$，则“若 $X$ 为 $x$ 时 $Y$ 会取的值”等于 $Y$ 的实际值。于是，结构进路把干预 $\doxp(x)$ 视为改变模型（与分布）但保持一切变量不变的操作，而潜在结局进路则把 $\doxp(x)$ 之下的变量 $Y$ 视为另一个变量 $Y(x)$，它通过 \textup{(3.52)} 之类的关系与 $Y$ 松散相连。第七章将利用 $Y(x,u)$ 的结构解释表明：在所有方面把反事实当作随机变量处理确实是正当的；而且诸如 \textup{(3.52)} 的一致性约束乃是从结构解释推出的定理，无须再考虑任何其他约束。

为交流实质性因果知识，潜在结局分析者必须把因果假设表述为对 $P^*$ 的约束，通常取涉及反事实变量之条件独立断言的形式。例如，为传达这样的理解——在依从性不完美的随机临床试验中（见图 3.7(b)），受试者对处理（$X$）的反应方式（$Y$）与处理分配（$Z$）统计独立——潜在结局分析者会写 $Y(x)\indep Z$。同样，为传达“分配是随机化的、因而与受试者如何依从分配无关”这一理解，潜在结局分析者会使用独立约束 $Z\indep X(z)$。

这一类约束的集合有时足以使所关心的查询有唯一解；在另一些情形中，则只能得到解的界。例如，若能合理地假定协变量集 $Z$ 满足条件独立
\begin{equation}
Y(x)\indep X\mid Z \tag{3.53}
\end{equation}
（Rosenbaum and Rubin 1983 称之为“条件可忽略性”），则可用 \textup{(3.52)} 容易地求出因果效应 $P^*(Y(x)=y)$，得到\footnote{Gibbard and Harper (1976, p.~157) 用“可忽略性假设”$Y(x)\indep X$ 导出了等式 $P(Y(x)=y)=P(y\mid x)$。}
\begin{equation}
\begin{aligned}
P^*(Y(x)=y) &= \sum_{z}P^*(Y(x)=y\mid z)P(z)\\
&= \sum_{z}P^*(Y(x)=y\mid x,z)P(z)\\
&= \sum_{z}P^*(Y=y\mid x,z)P(z)\\
&= \sum_{z}P(y\mid x,z)P(z).
\end{aligned} \tag{3.54}
\end{equation}
最后一式不含任何反事实量（因而可以去掉 $P^*$ 上的星号），并且恰与由后门准则得到的调整公式 \textup{(3.19)} 一致。然而，导出 \textup{(3.54)} 的关键——条件可忽略性假设（式 \textup{(3.53)}）——并不容易理解或确认。用实验的隐喻来转述，该假设读作：具有属性 $Z$ 的个体对处理 $X=x$ 会作出的反应，与该个体实际接受的处理无关。

\S 3.6.2 解释了为何这一进路可能吸引传统统计学家：不必为因果表达式构造新词汇与新逻辑，潜在结局框架中的一切数学运算都在概率演算的安全边界内进行。其缺点在于，它要求用反事实变量之间的独立关系来表达朴素的因果知识。当反事实变量不被视为某个更深层的、基于过程的模型的副产品时，就很难确认：一切相关的反事实独立判断是否都已被表述，\footnote{在图 3.7(b) 的例子中，一种典型的疏忽是写成 $Z\indep Y(x)$ 与 $Z\indep X(z)$，而不是按 \textup{(3.56)} 所规定的 $Z\indep\{Y(x),X(z)\}$。}已表述的判断是否冗余，或这些判断是否自洽。此类反事实判断的引出，可以借助下述从图出发的翻译而系统化（更多的关系见 \S 7.1.4）。

图在方程与概率函数 $P(u)$ 二者之中编码实质信息；前者编码为缺失的箭头，后者编码为缺失的虚线弧。因果图 $G$ 中每个子—父家族 $(\PA_i,X_i)$ 对应于 \textup{(3.50)} 中模型 $M$ 的一个方程。因此，缺失的箭头编码排除性假设，即断言：把被排除的变量加入某方程，不会改变该方程所描述之假设实验的结果。缺失的虚线弧编码两个或更多方程中扰动项之间的独立关系。例如，节点 $Y$ 与节点集 $\{Z_1,\dots,Z_k\}$ 之间没有虚线弧，意味着相应的背景变量 $U_Y$ 与 $\{U_{Z_1},\dots,U_{Z_k}\}$ 在 $P(u)$ 中独立。

这些假设可以用两条简单规则翻译为潜在结局记号（Pearl 1995a, p.~704）；第一条解释图中缺失的箭头，第二条解释缺失的虚线弧。
\begin{enumerate}[label=\arabic*.,leftmargin=1.8em]
\item \textbf{排除限制}：对每个具有父节点 $\PA_Y$ 的变量 $Y$，以及对与 $\PA_Y$ 不相交的任一变量集 $S$，有
\begin{equation}
Y(\pa_Y)=Y(\pa_Y,s). \tag{3.55}
\end{equation}
\item \textbf{独立限制}：若 $Z_1,\dots,Z_k$ 是任一不经虚线弧与 $Y$ 相连的节点集，则\footnote{这一限制实际上更强，它同时适用于 $\PA$ 变量的一切实例化。例如，$X\indep Y(\pa_Z)$ 应被理解为 $X\indep\{Y(\pa'_Z),Y(\pa''_Z),Y(\pa'''_Z),\dots\}$，其中 $(\pa'_Z),(\pa''_Z),(\pa'''_Z),\dots$ 是集合 $\PA_Z$ 可能取得的各个值。}
\begin{equation}
Y(\pa_Y)\indep\{Z_1(\pa_{Z_1}),\dots,Z_k(\pa_{Z_k})\}. \tag{3.56}
\end{equation}
\end{enumerate}
独立限制把 $U_Y$ 与 $\{U_{Z_1},\dots,U_{Z_k}\}$ 之间的独立，翻译为相应潜在结局变量之间的独立。这来自如下观察：一旦设定其父节点，$\{Y,Z_1,\dots,Z_k\}$ 中的变量便与其相应方程中的 $U$ 项处于函数关系。

例如，图 3.5 所示的模型具有如下父节点集：
\begin{equation}
\PA_X=\varnothing,\qquad \PA_Z=\{X\},\qquad \PA_Y=\{Z\}. \tag{3.57}
\end{equation}
因此，排除限制译为：
\begin{equation}
Z(x)=Z(y,x), \tag{3.58}
\end{equation}
\begin{equation}
X(y)=X(x,y)=X(z)=X, \tag{3.59}
\end{equation}
\begin{equation}
Y(z)=Y(z,x); \tag{3.60}
\end{equation}
而 $Z$ 与 $\{Y,X\}$ 之间没有虚线弧，则译为独立限制
\begin{equation}
Z(x)\indep\{Y(z),X\}. \tag{3.61}
\end{equation}
给定 $P^*$ 上足够多的此类限制，分析者便尝试用标准概率演算、连同把反事实变量与其可测对应物联系起来逻辑约束（如 \textup{(3.52)}），来计算因果效应 $P^*(Y(x)=y)$。这些约束可以作为公理或推理规则，用来把形如 $P^*(Y(x)=y)$ 的因果效应表达式变换为只涉及可测变量的表达式。一旦找到这样的变换，相应的因果效应便是可识别的，因为此时 $P^*$ 归结为 $P$。

于是自然产生这样的问题：潜在结局分析者所用的约束是否完备——即它们是否足以导出关于因果过程、干预与反事实的每一句有效陈述。要回答这个问题，就必须相对于更基本的数学对象（如可能世界，\S 1.4.4；或结构方程，式 \textup{(3.51)}）来定义反事实陈述的有效性。然而在标准的潜在结局框架中，完备性问题仍是悬而未决的，因为 $Y(x,u)$ 被当作原初概念，而诸如 \textup{(3.52)} 的一致性约束——尽管对英语表达“若 $X$ 曾为 $x$”而言似乎合理——并未从更深的数学对象中导出。第七章将解决这一完备性问题：那里从 \textup{(3.51)} 赋予 $Y(x,u)$ 的结构语义出发，导出了一组充分必要的公理。

在评估结构方程与潜在结局模型的历史发展时，结构方程相对于潜在结局模型所提供的概念清晰性，其重要性无论怎样强调都不过分。读者若想在某个熟悉情境中判断 \textup{(3.61)} 的条件是否成立，便可体会这一点。该条件读作：“若 $X$ 曾为 $x$，$Z$ 会取得的值，与 $X$ 以及‘若 $Z$ 曾为 $x$，$Y$ 会取得的值’二者联合独立。”（在结构表示中，这句话读作：“除 $X$ 本身之外，$Z$ 与 $X$ 或 $Y$ 都不共享原因，如图 3.5 所示。”）一想到必须表述、辩护并管理这类令人生畏的反事实关系，或许就能解释：为何因果推断这一事业在普通流行病学家与统计学家眼中显得如此可畏而令人绝望；也能解释为何多数经济学家与社会科学研究者继续使用结构方程，而不用 Holland (1988)、Angrist 等 (1996) 与 Sobel (1998) 所倡导的潜在结局替代方案。另一方面，一旦问题被恰当形式化，潜在结局记号所提供的代数机械在精炼假设、导出反事实概率、验证结论是否由前提推出等方面可以相当强大——正如我们将在第九章演示的那样。\textup{(3.51)}—\textup{(3.56)} 所给的翻译，是统一两个阵营的关键，应有助于研究者结合两种进路的最佳特性。

\subsection{与 Robins 的 G-估计的关系}

在潜在结局框架中所作的研究中，在精神上最接近本章所述结构分析的，是 Robins 关于“因果解释结构化树图”的工作（Robins 1986, 1987）。Robins 第一个认识到 Neyman 的反事实记号 $Y(x)$ 作为一般因果推断数学语言的潜力，并用它把 Rubin (1978) 的“时间无关处理”模型推广到含直接效应与间接效应、时变处理、伴随变量与结局的研究。

Robins 考虑一个按时间排序的离散随机变量集 $V=\{V_1,\dots,V_M\}$（如图 3.3），并追问：在什么条件下可以识别控制策略 $g:X=x$ 对结局 $Y\subseteq V\backslash X$ 的效应，其中 $X=\{X_1,\dots,X_K\}\subseteq V$ 是所关心的、按时间排序且可能可操纵的处理变量。$X=x$ 对 $Y$ 的因果效应被表述为概率
\[
P(y\mid g=x)\triangleq P\{Y(x)=y\},
\]
其中反事实变量 $Y(x)$ 表示：若处理变量 $X$ 曾为 $x$，结局变量 $Y$ 会取的值。

Robins 证明：若 $X$ 的每个分量 $X_k$ 都是“在给定过去时随机分配的”，则 $P(y\mid g=x)$ 可由分布 $P(v)$ 识别。这一概念阐述如下。设 $L_k$ 为出现在 $X_{k-1}$ 与 $X_k$ 之间的变量，$L_1$ 为 $X_1$ 之前的变量。记 $\bar{L}_k=(L_1,\dots,L_k)$、$\bar{L}=\bar{L}_K$、$\bar{X}_k=(X_1,\dots,X_k)$，并令 $\bar{L}_0$、$\bar{X}_0$、$V_0$ 恒为零。若下述关系成立，则称处理 $X_k=x_k$ 是给定过去时随机分配的：
\begin{equation}
(Y(x)\indep X_k\mid \bar{L}_k,\bar{X}_{k-1}=x_{k-1}). \tag{3.62}
\end{equation}
Robins 进一步证明：若 \textup{(3.62)} 对每个 $k$ 都成立，则因果效应由
\begin{equation}
P(y\mid g=x)=\sum_{\bar{l}_K}P(y\mid \bar{l}_K,x_K)\prod_{k=1}^{K}P(l_k\mid \bar{l}_{k-1},x_{k-1}) \tag{3.63}
\end{equation}
给出，他称此式为“G-计算算法公式”。该式可通过迭代施用条件 \textup{(3.62)} 导出，如同导出 \textup{(3.54)} 那样。若 $X$ 为单变量，则 \textup{(3.63)} 化简为标准调整公式
\[
P(y\mid g=x)=\sum_{l_1}P(y\mid x,l_1)P(l_1),
\]
与 \textup{(3.54)} 平行。同样，在图 3.3 的特别结构中，\textup{(3.63)} 化简为 \textup{(3.18)}。

为把这一结果置于本章分析的语境中，我们需要聚焦于条件 \textup{(3.62)}——它促成了 Robins 对 \textup{(3.63)} 的推导——并追问：这一形式的反事实独立性能否获得有意义的图形解释。答案将在第四章（定理 4.4.1）给出，那里我们导出了识别一个\emph{计划}（即一序列行动）之效应的图形条件。该条件读作：若从 $X_k$ 到 $Y$ 的每条“避开行动的”后门路径都被 $X_k$ 之非后代的某个子集 $L_k$ 阻断，则 $P(y\mid g=x)$ 可识别并由 \textup{(3.63)} 给出。（所谓“避开行动”，我们指路径中不含任何进入晚于 $X_k$ 的 $X$ 变量的箭头。）第十一章（\S 11.4.2）用例子表明，这一“序贯后门”准则比 \textup{(3.62)} 所给者更一般。

本章引入的结构分析从一个新的理论视角支持并推广了 Robins 的结果。第一，在技术层面，这一分析提供了系统的方式来处理 Robins 的出发点 \textup{(3.62)} 不适用的模型，例如图 3.8(d)—(g)。第二，在概念层面，结构框架代表着从反事实独立性的词汇，向过程与机制的词汇的根本转变——人类知识正是以这种词汇编码的。前者要求人们确认诸如 \textup{(3.62)} 这样晦涩的关系，后者则用缺失连线的生动图形语言来表达同样的关系。Robins 的开创性研究表明，为恰当处理时变处理的多阶段问题，晦涩的“可忽略性”条件 \textup{(3.53)} 应当被分解为其序贯成分。这导致了定理 4.4.5 的序贯后门准则。

\section*{个人感想与致谢}
\addcontentsline{toc}{section}{个人感想与致谢}

本章所述的工作，源于两个彻底改变了我对因果性之态度的简单想法。第一个想法出现在 1990 年夏天，当时我正与 Tom Verma 合作研究“推断因果性的理论”（Pearl and Verma 1991；另见第二章）。我们尝试用函数对应物 $x_i=f_i(\pa_i,u_i)$ 替换父子关系 $P(x_i\mid\pa_i)$，忽然间一切各就各位：我们终于有了一个数学对象，可以把物理机制的熟悉性质赋予它，而不必再依赖那些在贝叶斯网络研究中与之打交道已久的、滑溜的认识论概率 $P(x_i\mid\pa_i)$。当时正在写博士论文的 Danny Geiger 惊讶地问：“确定性方程？真的确定性的？”尽管我们知道确定性结构方程在计量经济学中有悠久历史，我们却把这种表示视为过去的遗物。对 1990 年代初的 UCLA 而言，把贝叶斯网络的语义放在确定性基础之上，似乎是最好意义上的异端。

第二个简单想法来自 Peter Spirtes 在国际科学哲学大会（瑞典乌普萨拉，1991）上的演讲。在一张幻灯片中，Peter 演示了当一个变量被操纵时因果图会如何变化。对我来说，Spirtes 的那张幻灯片——与确定性结构方程相结合——就是展开操纵论因果说明的钥匙，并引出了本章所描述的大部分探索。

我还应当提到另一件促成本章的事。1993 年初，我读了 Arthur Goldberger 与 Nanny Wermuth 关于结构方程含义的激烈辩论（Goldberger 1992; Wermuth 1992）。我忽然领悟到，经济学家与统计学家之间长达一个世纪的紧张关系，源于一个简单的语义混淆：统计学家把结构方程读作关于 $E(Y\mid x)$ 的陈述，而经济学家把它读作 $E(Y\mid\doxp(x))$。这就解释了为什么统计学家声称结构方程毫无意义，而经济学家反驳说统计学没有实质内容。我写了一份技术报告《论结构方程的统计解释》（Pearl 1993c），希望看到两个阵营和解相拥。什么也没发生。争论中的统计学家继续坚持：凡不被解释为 $E(Y\mid x)$ 的东西就毫无意义。相反，经济学家仍在琢磨：他们一直想说的，是否就是 $\doxp(x)$。

给予鼓励的同行所获得的正式承认实在太少，考虑到他们对被鼓励者的巨大影响。我必须借此机会感谢四位同行，他们在 $\doxp(x)$ 算子流行之前就看到了其中的清晰性：Steffen Lauritzen、David Freedman、James Robins 与 Philip Dawid。Phil 尤其有勇气把我的论文刊于 \textit{Biometrika}（Pearl 1995a）——那份由因果性最凶猛的对手 Karl Pearson 创办的期刊。

\section*{第二版后记}
\addcontentsline{toc}{section}{第二版后记}

\subsection*{完整的识别结果}

一个推广了本章所建立之一切准则的关键识别条件，由 Jin Tian 导出。它读作：

\begin{theorem}[Tian and Pearl 2002a]
识别因果效应 $P(y\mid\doxp(x))$ 的一个充分条件是：$X$ 与其任一子节点之间不存在双向路径（即完全由双向弧构成的路径）。\footnote{施用此准则之前，可以先从因果图中删去一切不是 $Y$ 之祖先的节点。}
\end{theorem}

值得注意的是，该定理断言：只要 $X$ 的每个子节点（在通往 $Y$ 的路径上）都不能由 $X$ 经双向路径到达，那么无论图多么复杂，因果效应 $P(y\mid\doxp(x))$ 都可识别。本章讨论的一切识别准则都是该定理所定义准则的特别情形。例如在图 3.5 中，$P(y\mid\doxp(x))$ 可识别，因为从 $X$ 到 $Z$（$X$ 的唯一子节点）的两条路径都不是双向的。相反，在图 3.7 中，存在一条从 $X$ 到 $Z_1$、只经过双向弧的路径，因而违反了定理 3.6.1 的条件，$P(y\mid\doxp(x))$ 不可识别。

注意图 3.8 中的所有图都满足上述条件，而图 3.9 中没有一个满足。Tian and Pearl (2002a) 进一步证明：当 $V$ 包含除 $X$ 之外的一切变量时，该条件对 $P(y\mid\doxp(x))$ 的识别既是充分的也是必要的。对 $P(w\mid\doxp(z))$（$W$、$Z$ 为任意两个集合）的识别的充分必要条件，由 Shpitser and Pearl (2006b) 建立。随后，一个完整的图形准则被建立起来，用以判定条件干预分布——即形如 $P(y\mid\doxp(x),z)$（$X,Y,Z$ 为任意变量集）的表达式——的可识别性（Shpitser and Pearl 2006a）。

这些结果构成了图模型中因果效应的完整刻画。它们为我们提供了多项式时间算法，用以判定任意诉诸 $\doxp(x)$ 算子的量在给定半马尔可夫模型中是否被识别，若被识别，其估计量又是什么。值得注意的是，这些结果的一个推论还表明：do 演算是完备的，即量 $Q=P(y\mid\doxp(x),z)$ 可识别，当且仅当它可以用定理 3.4.1 的三条规则化简为无 do 的表达式。\footnote{这一点由 Huang and Valtorta (2006) 独立确立。}Tian and Shpitser (2010) 对这些结果作了全面总结。

\subsection*{应用与批评者}

本章所发展概念的浅近介绍见 Pearl (2003c) 与 Pearl (2009a)。因果图在流行病学中的应用见 Robins (2001)、Hernán 等 (2002)、Hernán 等 (2004)、Greenland and Brumback (2002)、Greenland 等 (1999a,b)、Kaufman 等 (2005)、Petersen 等 (2006)、Hernández-Díaz 等 (2006)、VanderWeele and Robins (2007) 以及 Glymour and Greenland (2008)。

前门准则（\S 3.3.2）的有趣应用，见于社会科学（Morgan and Winship 2007）与经济学（Chalak and White 2006）。

“潜在结局”进路的一些倡导者最不愿接受把图或结构方程作为因果分析的基础；由于缺乏这些概念工具，他们无法处理协变量选择问题（Rosenbaum 2002, p.~76; Rubin 2007, 2008a），并进而把重要的科学概念斥为“定义不清”“欺骗性的”“令人困惑的”（Holland 2001; Rubin 2004, 2008b），乃至更甚（Rubin 2009）。Lauritzen (2004) 与 Heckman (2005) 批评了这种态度；Pearl (2009a,b, 2010a) 阐明了其有害后果。

同样令人费解的是，某些哲学家（Cartwright 2007; Woodward 2003）与经济学家（Heckman 2005）担心 do 算子过于局部，无法刻画复杂的现实政策干预——这类干预有时同时影响若干机制，且常涉及条件决策、不完美控制与多重行动。这些担忧来自把关系的数学定义（如因果效应）与该关系在物理世界中可检验的技术可行性混为一谈。do 算子确实是一件理想的数学工具（与微分学中的导数不无相似），然而它使我们能够指定并分析极其复杂的干预策略。读者将在第四章找到此类策略的例子，对这一问题的进一步讨论见第十一章（\S 11.4.3—11.4.6 与 \S 11.5.4）。

\section*{本章主要结果导览}
\addcontentsline{toc}{section}{本章主要结果导览}

本章的三个关键结果是：1. 混杂的控制；2. 政策的评估；3. 反事实的评估。
\begin{enumerate}[label=\arabic*.,leftmargin=1.8em]
\item 控制混杂偏倚的问题，通过后门条件（定理 3.3.2，第 79—80 页）得到解决——这是一个选择协变量集 $Z$ 的准则，若对之作调整，便可获得因果效应的无偏估计。
\item 政策评估问题——从非实验数据预测干预的效果——通过 do 演算（定理 3.4.1，第 85—86 页）及其所蕴含的图形准则（定理 3.3.4，第 83 页；定理 3.6.1，第 105 页）得到解决。do 演算的完备性意味着：任何（非参数）政策评估问题，若不为某个可识别图或等价的因果假设集所支持，便可被证明为“不可解”。
\item 最后，式 \textup{(3.51)} 为反事实提供了形式语义，通过它，反事实的联合概率可以在科学理论的框架内被定义与求值（见第七章）。这一语义将使我们能够发展一系列反事实分析技术（第八至十一章），包括中介公式（式 \textup{(4.17)}—\textup{(4.18)}）——在非线性模型中评估因果通路的关键工具。
\end{enumerate}
